% Bộ đề HSA mẫu nguyên bản cho renderer/UI.
% Nội dung dưới đây được viết mới theo các dạng quan sát được; không sao chép
% câu hỏi, đáp án hay response của đề tham chiếu.
% Có thể mở rộng mỗi phần lên số lượng chính thức trong workspace LaTeX.

\phan1
\begin{ex}
Một cửa hàng ghi nhận số lượt khách trong bốn ngày lần lượt là $42$, $55$, $48$ và $65$. Số lượt khách trung bình trong bốn ngày là
\choice
{$50$}
{\True $52{,}5$}
{$54$}
{$55$}
\loigiai{Cộng bốn số liệu rồi chia cho $4$.}
\end{ex}

\begin{ex}
Cho bảng số liệu về lượng nước sử dụng (lít) của một gia đình:
\begin{center}
\begin{tabular}{|c|c|c|c|}
\hline
Ngày & Thứ hai & Thứ ba & Thứ tư \\ \hline
Lượng nước & $120$ & $135$ & $105$ \\ \hline
\end{tabular}
\end{center}
Nếu mục tiêu mỗi ngày là $110$ lít thì ngày nào vượt mục tiêu nhiều nhất?
\choice
{Thứ hai}
{\True Thứ ba}
{Thứ tư}
{Cả ba ngày như nhau}
\end{ex}

\begin{ex}
Một cảm biến ghi nhiệt độ theo công thức $T(t)=24+3\sin\dfrac{\pi t}{12}$, với $t$ tính bằng giờ. Nhiệt độ cao nhất trong chu kì là bao nhiêu độ C?
\choice
{$21$}
{$24$}
{\True $27$}
{$30$}
\end{ex}

\begin{ex}
Một xe đạp đi quãng đường $18$ km trong $45$ phút. Vận tốc trung bình của xe là bao nhiêu km/h?
\choice
{$18$}
{\True $24$}
{$27$}
{$30$}
\end{ex}

\begin{ex}
Một hộp có $5$ thẻ màu xanh và $3$ thẻ màu vàng. Lấy ngẫu nhiên một thẻ, xác suất lấy được thẻ vàng là
\choice
{$\dfrac{3}{5}$}
{\True $\dfrac{3}{8}$}
{$\dfrac{5}{8}$}
{$\dfrac{1}{3}$}
\end{ex}

\begin{ex}
Một bể nước có $240$ lít. Mỗi phút máy bơm thêm $12$ lít và van xả bớt $4$ lít. Sau bao nhiêu phút bể có $400$ lít?
\cham{20}
\loigiai{Mỗi phút lượng nước tăng $8$ lít nên cần $(400-240)/8=20$ phút.}
\end{ex}

\begin{ex}
Một hình chữ nhật có chiều dài $14$ cm và chiều rộng $9$ cm. Chu vi của hình chữ nhật là bao nhiêu cm?
\cham{46}
\end{ex}

\begin{ex}
Một khoản tiền $12$ triệu đồng được chia đều cho $8$ nhóm. Mỗi nhóm nhận bao nhiêu triệu đồng?
\cham{1,5}
\end{ex}

\phan2
\begin{ex}
Chọn từ có nghĩa khác nhóm với các từ còn lại.
\choice
{\True Kiên trì}
{Bền bỉ}
{Nhẫn nại}
{Dai dẳng}
\end{ex}

\begin{ex}
Chọn cụm từ có cấu tạo khác các cụm từ còn lại.
\choice
{Niềm tin}
{Màu xanh}
{\True Rất nhanh}
{Lòng biết ơn}
\end{ex}

\begin{ex}
Chọn cặp từ thích hợp nhất điền vào chỗ trống: Một quyết định tốt cần dựa trên dữ liệu đáng tin cậy và một quá trình ________ trước khi hành động.
\choice
{cảm tính / vội vàng}
{\True phân tích / cân nhắc}
{ngẫu nhiên / im lặng}
{thụ động / bỏ qua}
\end{ex}

\begin{ex}
Xác định từ dùng sai trong câu: “Những kết quả mới đã mở ra một hướng tiếp cận đầy triển vọng cho nhóm nghiên cứu.”
\choice
{kết quả}
{mới}
{\True đầy đặn}
{triển vọng}
\end{ex}

\begin{baidoc}
\item[] Dựa vào thông tin dưới đây và trả lời các câu hỏi từ câu {start} - {end}:
\tenbaidoc{Khu vườn trên mái}
Một khu vườn trên mái có thể làm giảm nhiệt độ bề mặt của tòa nhà, giữ lại một phần nước mưa và tạo thêm nơi sống cho côn trùng thụ phấn. Hiệu quả của khu vườn phụ thuộc vào độ sâu của lớp đất, loại cây được trồng và cách chăm sóc. Nếu thiết kế không phù hợp, mái nhà có thể chịu tải quá lớn hoặc hệ thống thoát nước bị tắc.
\begin{ex}
Mục đích chính của văn bản là gì?
\choice
{Kể lại lịch sử của các tòa nhà cao tầng}
{\True Giới thiệu lợi ích và điều kiện để vận hành vườn trên mái}
{Chứng minh mọi mái nhà đều thích hợp để trồng cây}
{Cảnh báo không nên sử dụng nước mưa}
\end{ex}
\begin{ex}
Yếu tố nào cần được xem xét để tránh mái nhà chịu tải quá lớn?
\choice
{Màu sắc của hoa}
{Số lượng côn trùng}
{\True Độ sâu của lớp đất và khối lượng cây}
{Chiều cao của tòa nhà}
\end{ex}
\end{baidoc}

\begin{ex}
Đọc khổ thơ sau và xác định nội dung chính:
\begin{center}
Sớm mai trên phố nhiều cây\\
Một làn gió mỏng chạm vai người\\
Ai đi giữ lại màu xanh ấy\\
Cho những ngày dài bớt nắng thôi.
\end{center}
\choice
{Nỗi buồn trước sự đổi thay của phố}
{\True Mong muốn gìn giữ không gian xanh trong đời sống}
{Niềm vui của người đi du lịch}
{Lời kể về một buổi chợ đông}
\end{ex}

\phan3
% Tổ hợp Khoa học: các câu mẫu dưới đây đại diện cho Vật lí, Hóa học,
% Sinh học, Lịch sử và Địa lí; workspace có thể chọn đúng ba chủ đề.
\begin{ex}
Một vật khối lượng $2$ kg chuyển động với gia tốc $3$ m/s$^2$. Độ lớn hợp lực tác dụng lên vật là
\choice
{$1{,}5$ N}
{$5$ N}
{\True $6$ N}
{$9$ N}
\end{ex}

\begin{ex}
Một bình kín chứa khí ở nhiệt độ không đổi. Khi thể tích giảm còn một nửa thì áp suất của khí
\choice
{giảm còn một nửa}
{không đổi}
{\True tăng gấp đôi}
{tăng gấp bốn}
\end{ex}

\begin{ex}
Phản ứng giữa dung dịch acid và base tạo ra muối và nước được gọi là phản ứng
\choice
{oxi hóa - khử}
{phân hủy}
{\True trung hòa}
{trùng hợp}
\end{ex}

\begin{ex}
Trong tế bào, bào quan thực hiện phần lớn quá trình hô hấp hiếu khí là
\choice
{ribosome}
{lục lạp}
{\True ti thể}
{không bào}
\end{ex}

\begin{ex}
Đánh dấu Đúng hoặc Sai cho các nhận định về quang hợp.
\choiceTF
{\True Quang hợp sử dụng năng lượng ánh sáng.}
{\True Sản phẩm hữu cơ có thể được tích lũy dưới dạng tinh bột.}
{Quang hợp chỉ xảy ra ở tế bào động vật.}
{\True Diệp lục hấp thụ một phần năng lượng ánh sáng.}
\end{ex}

\begin{ex}
Ghép mỗi khái niệm ở cột trái với mô tả phù hợp ở cột phải.
\match
{Quang hợp}
{Hô hấp tế bào}
{Phản xạ ánh sáng}
\with
[2]{Giải phóng năng lượng từ chất hữu cơ trong tế bào.}
[3]{Ánh sáng đổi hướng khi gặp bề mặt phản xạ.}
[1]{Tạo chất hữu cơ nhờ năng lượng ánh sáng.}
\end{ex}

% Tổ hợp Tiếng Anh có thể dùng cùng renderer bằng cách thay phần 3 trong
% workspace. Các dạng mẫu: grammar, dialogue, sentence paraphrase, reading.
\begin{ex}
Choose A, B, C or D to complete the sentence: The research team ________ the results before publishing them.
\choice
{check}
{checking}
{\True checked}
{has check}
\end{ex}

\begin{ex}
Choose the closest meaning to the sentence: “The project was delayed because the data had not been verified.”
\choice
{The project verified the data before it started.}
{\True The lack of verified data caused the project to be late.}
{The data arrived after the project was completed.}
{The project did not use any data.}
\end{ex}
